\documentclass[11pt,a4paper]{article}

% ============================================================
% Packages (core only — texlive-latex-base + texlive-base)
% ============================================================
\usepackage{amsmath, amssymb, amsthm}
\usepackage[a4paper, margin=2.5cm]{geometry}
\usepackage{array}

% ============================================================
% Theorem environments
% ============================================================
\theoremstyle{definition}
\newtheorem{definition}{Definition}[section]
\newtheorem{remark}{Remark}[section]
\newtheorem{proposition}{Proposition}[section]

% ============================================================
% Notation macros
% ============================================================
\newcommand{\R}{\mathbb{R}}
\newcommand{\E}{\mathbb{E}}
\newcommand{\N}{\mathbb{N}}
\newcommand{\norm}[1]{\left\lVert #1 \right\rVert}
\newcommand{\inner}[2]{\langle #1,\, #2 \rangle}
\newcommand{\corr}{\operatorname{corr}}
\newcommand{\sgn}{\operatorname{sign}}
\newcommand{\code}[1]{\texttt{\detokenize{#1}}}

% ============================================================
\begin{document}

\title{\textbf{Perception Pressure and Resonance}\\[4pt]
\large Decomposable, Trust-Weighted Measures of Media Narrative\\
from a Multi-Perspective News-Aggregation System}

\author{Yigit Onat\\
\textit{Independent Researcher}\\
\texttt{yionat@gmail.com}}
\date{June 2026}
\maketitle

% ============================================================
\begin{abstract}
We introduce two quantitative measures of media narrative computed continuously from a live,
multi-source news stream, and the system---\emph{Prism}---that produces them. \emph{Resonance} scores
the media impact of a single story by combining source diversity, source credibility, audience
engagement, and recency. \emph{Perception Pressure} is a per-concept longitudinal signal,
$P(K,t) = \text{salience}(K,t)\times\text{valence}(K,t)$, that couples \emph{how much} a concept is
covered (a non-negative attention term) with \emph{how} it is framed (a credibility-weighted
directional sentiment term in $[-1,+1]$). The design has three properties we argue are novel in
combination: it multiplicatively couples attention and directional sentiment, so the signal is
\emph{signed} and its magnitude is meaningful; source diversity enters as a \emph{multiplier}, so a
single outlet cannot inflate the measure (multi-perspective by construction); and every component is
persisted independently, so the measure is fully decomposable and auditable rather than a black box.
Credibility is not hand-assigned but \emph{earned} through an automated cross-corroboration
lifecycle. The system is realized as a five-agent pipeline (discover, analyze, track, personalize,
deliver) coordinated by a database state machine. We position Perception Pressure against
volume-only attention proxies and direction-only sentiment indices, formalize its construction, and
specify a downstream application: treating $P(K,t)$ as an \emph{exogenous narrative signal} whose
incremental predictive value for financial returns can be measured under a leakage-controlled
information-theoretic protocol. We are explicit that this is a methods-and-system paper: the
instrument is fully specified and parameterized for ablation, and empirical validation under the
stated protocol is the next step rather than a claim made here.
\end{abstract}

\smallskip
\noindent\textbf{Keywords:} media narrative, attention, sentiment, agenda-setting, source
credibility, news analytics, exogenous signal, computational social science

\smallskip
\noindent\textbf{JEL Classification:} C43, D83, G14, L82

\tableofcontents
\bigskip

% ============================================================
\section{Introduction}
\label{sec:intro}
% ============================================================

\subsection{Measuring narrative, not just volume or tone}

Two literatures measure the media's relationship to attention and belief, and each captures only
half of the object of interest. \emph{Attention} proxies---search-volume indices, mention counts---
measure \emph{how much} a topic is discussed but are sign-blind: a surge of coverage looks identical
whether the framing is favorable or hostile. \emph{Sentiment} indices measure \emph{how} a topic is
framed but are typically volume-blind and treat all sources alike, so a single fringe outlet and a
wire service count equally. What is missing is a single measure that is simultaneously
\emph{directional} (positive vs.\ negative framing), \emph{magnitude-meaningful} (broad, authoritative,
fresh coverage weighs more than thin coverage), and \emph{credibility-aware} (trusted corroboration
counts for more than isolated amplification).

This paper introduces such a measure, \emph{Perception Pressure}, and the system that computes it.
The construction is deliberately simple: it multiplies a non-negative \emph{salience} term (weighted
attention) by a \emph{valence} term (trust-weighted mean sentiment) bounded in $[-1,+1]$. The product
is signed, its sign is the net framing of the concept, and its magnitude grows with the breadth,
authority, and freshness of coverage. Because it is a longitudinal signal sampled on a fixed cadence,
its time-derivative---\emph{momentum}---distinguishes a narrative building from a backlash forming.

\subsection{Contributions}

\begin{enumerate}
  \item \textbf{Perception Pressure}, a per-concept, signed, magnitude-meaningful narrative signal
        $P(K,t) = \text{salience}\times\text{valence}$, with a credibility-weighted valence and an
        explicit momentum (Section~\ref{sec:perception}).
  \item \textbf{Resonance}, a per-story media-impact score in which source diversity is a
        \emph{multiplier}, operationalizing ``multi-perspective by construction''
        (Section~\ref{sec:resonance}).
  \item \textbf{Earned credibility}: an automated source-trust lifecycle in which a source's weight
        emerges from cross-source corroboration rather than being assigned a priori
        (Section~\ref{sec:trust}).
  \item \textbf{Decomposability as a design constraint}: every sub-component (salience, valence,
        breadth, decay, trust) is persisted independently, so the measure is auditable rather than a
        black box (Section~\ref{sec:properties}).
  \item \textbf{A downstream application and its evaluation protocol}: treating $P(K,t)$ as an
        exogenous signal whose incremental predictive value for markets is measured under a
        leakage-controlled information-theoretic protocol (Sections~\ref{sec:application}--\ref{sec:eval}).
\end{enumerate}

This is a methods-and-system paper. The metrics are implemented and every parameter is a tunable
knob suitable for ablation; we report no empirical performance figures, and Section~\ref{sec:eval}
specifies the evaluation that would substantiate the measure's utility.

% ============================================================
\section{Related Work}
\label{sec:related}
% ============================================================

\paragraph{Agenda-setting and attention.}
McCombs and Shaw (1972) established that media coverage shapes which issues the public deems
important. Da, Engelberg, and Gao (2011) operationalize \emph{investor} attention with search
volume; Preis, Moat, and Stanley (2013) and Choi and Varian (2012) use Google Trends to nowcast
behavior. These are volume measures: they capture salience but not framing.

\paragraph{Media tone and markets.}
Tetlock (2007) shows that media \emph{pessimism} forecasts returns; Antweiler and Frank (2004) and
Bollen, Mao, and Zeng (2011) study message-board and social tone; Loughran and McDonald (2011)
provide finance-specific sentiment lexicons. These are direction measures, typically source-agnostic
and not coupled to attention. Perception Pressure differs by \emph{multiplying} attention and tone
and by weighting tone by earned source credibility.

\paragraph{Text-based indices and slant.}
Baker, Bloom, and Davis (2016) build a policy-uncertainty index from newspaper text; Gentzkow,
Shapiro, and Taddy (2019) measure partisan slant from text. Our valence term is in this tradition but
is computed online, per concept, with a credibility weight, and is paired with a salience term into a
single signed signal.

\paragraph{Information-theoretic signal evaluation.}
Our downstream evaluation uses mutual information (Kraskov, St\"ogbauer, and Grassberger, 2004) and
false-discovery-rate control (Benjamini and Hochberg, 1995) to test predictive content without
overfitting---methods imported from the consuming system rather than invented here.

% ============================================================
\section{System Overview}
\label{sec:system}
% ============================================================

Prism is an AI multi-perspective news-aggregation platform whose stated principle is to make
editorial bias \emph{visible} rather than hidden: it performs no original reporting, attributes every
claim to a source, and presents multiple framings of each story side by side. Quantitatively, it
instruments the two narrative measures of this paper. It is realized as a five-agent pipeline,
coordinated not by a message queue but by a database state machine in which each story advances
through states \textsc{raw} $\to$ \textsc{analyzed} $\to$ \textsc{delivered}.

\begin{center}\small
\renewcommand{\arraystretch}{1.25}
\setlength{\tabcolsep}{5pt}
\begin{tabular}{>{\raggedright\arraybackslash}p{1.9cm} >{\raggedright\arraybackslash}p{6.6cm} >{\raggedright\arraybackslash}p{5.8cm}}
\hline
\textbf{Agent} & \textbf{Role} & \textbf{Method} \\
\hline
Discovery & pull and cluster articles by event; grow the source registry &
  news/RSS ingestion; lexical dedup (Jaccard, TF--IDF, entity overlap) \\
Analysis & summarize a cluster; extract per-source perspective, sentiment, bias, claims &
  large-language-model structured extraction with a quality score \\
Tracking & scan clusters for tracked concepts; compute the narrative measures &
  keyword matching; Resonance and Perception Pressure on a fixed cadence \\
Personal\-ization & rank stories per reader &
  explicit weighted scoring (interest, recency, diversity, resonance) \\
Delivery & compose and deliver an attributed briefing &
  templated generation; every claim links to its source \\
\hline
\end{tabular}
\end{center}

The tracking agent is the locus of this paper: on a fixed interval it recomputes, for each tracked
concept, the salience, valence, perception, and momentum series defined below and persists them as a
time-series. The other agents supply its inputs (clusters, perspectives, sentiments, source trust)
and consume its outputs (resonance feeds personalization).

% ============================================================
\section{Source Trust as an Earned Weight}
\label{sec:trust}
% ============================================================

Both measures weight evidence by source credibility. Crucially, credibility is not a static per-outlet
constant; it is earned through an auditable lifecycle:
\[
  \textsc{seed} \longrightarrow \textsc{candidate} \longrightarrow \textsc{probation}
  \longrightarrow \textsc{trusted} \;/\; \textsc{rejected}.
\]
A small registry of vetted outlets seeds the system. New domains observed in coverage become
candidates and, after repeated sightings, enter probation with a low initial trust. During probation
a source's articles are \emph{cross-validated}: an article is corroborated if it appears in a story
cluster that also contains a trusted source, and flagged if it is effectively single-sourced. Trust
during probation rises with the corroboration rate,
\begin{equation}
  \mathrm{trust} = 0.1 + 0.4\cdot\frac{\#\,\text{validated}}{\#\,\text{total}},
  \label{eq:probation}
\end{equation}
and after a probation window a source is promoted to \textsc{trusted} only if it has accumulated
enough corroborated articles at a sufficient rate; persistent contradiction demotes it. A bias label
is then inferred from the average sentiment of its perspectives. This makes the trust weight in the
measures below a function of demonstrated cross-source agreement rather than an editorial prior, and
it is the mechanism that makes the measures robust to single-source amplification.

% ============================================================
\section{Resonance: Per-Story Media Impact}
\label{sec:resonance}
% ============================================================

\begin{definition}[Resonance]
\label{def:resonance}
For a story cluster $c$ comprising mentions $m$ (one per article), with mention source $s_m$,
audience reaction count $\mathrm{re}_m$, and age $a_m$,
\begin{equation}
  \mathrm{Res}(c,t) = \underbrace{\mathrm{breadth}(c)}_{\text{diversity multiplier}}
       \cdot \sum_{m\in c} \mathrm{trust}(s_m)\,\cdot\,\mathrm{eng}(m)\,\cdot\,\mathrm{dec}(a_m),
  \label{eq:resonance}
\end{equation}
with
\begin{equation}
  \mathrm{breadth}(c) = \log_2\!\big(1 + n_{\text{sources}}(c)\big),\quad
  \mathrm{eng}(m) = 1 + \log_2\!\Big(1 + \tfrac{\mathrm{re}_m}{\text{median}}\Big),\quad
  \mathrm{dec}(a) = e^{-\lambda a},\;\; \lambda = \tfrac{\ln 2}{h},
  \label{eq:resonance-terms}
\end{equation}
where $n_{\text{sources}}(c)$ is the number of \emph{distinct} sources, $\mathrm{re}_m/\text{median}$
normalizes reactions to a platform median, and $h$ is a half-life (default $24$\,h).
\end{definition}

Three design choices follow from the platform's principle. \emph{Breadth as a multiplier}: a story
carried by one outlet, however loud, has $\mathrm{breadth} = \log_2 2 = 1$ and cannot resonate;
diversity is a precondition, not an additive bonus. \emph{Logarithmic engagement}: viral outliers are
compressed, so a single spike cannot dominate. \emph{Exponential decay}: impact is recency-weighted.
Derived quantities---\emph{momentum} $\mathrm{Res}(t)-\mathrm{Res}(t-\Delta t)$, peak resonance, and
persistence---track a story's life cycle. Every factor in \eqref{eq:resonance} is stored
independently, so resonance is queryable down to its components.

% ============================================================
\section{Perception Pressure: Per-Concept Narrative}
\label{sec:perception}
% ============================================================

Resonance scores a single story. Perception Pressure aggregates across all stories touching a tracked
concept $K$ into a longitudinal, signed signal.

\begin{definition}[Salience, valence, perception]
\label{def:perception}
For concept $K$ at time $t$, summing over clusters $c$ mentioning $K$ (with cluster breadth and decay
as above) and over each cluster's per-source perspectives $p$ (with trust $\mathrm{trust}(p)$ and
sentiment $\mathrm{sen}(p)\in[-1,1]$):
\begin{align}
  \text{salience:}\quad & A(K,t) = \sum_{c} \mathrm{breadth}(c)\,\mathrm{dec}(a_c) \;\;\ge 0,
    \label{eq:salience}\\[2pt]
  \text{valence:}\quad & V(K,t) =
    \frac{\sum_{c} \mathrm{breadth}(c)\,\mathrm{dec}(a_c)\,\sum_{p}\mathrm{trust}(p)\,\mathrm{sen}(p)}
         {\sum_{c} \mathrm{breadth}(c)\,\mathrm{dec}(a_c)\,\sum_{p}\mathrm{trust}(p)}
    \;\in [-1,1],
    \label{eq:valence}\\[2pt]
  \text{perception:}\quad & P(K,t) = A(K,t)\cdot V(K,t),
    \label{eq:perception}
\end{align}
with momentum $\dot P(K,t) = P(K,t) - P(K,t-1)$ over consecutive snapshots.
\end{definition}

The interpretation is direct. The sign of $P$ is the net framing of the concept; its magnitude grows
with strong, broad, authoritative, fresh coverage. $|P|\approx 0$ is genuinely ambiguous between
\emph{uncovered} (low salience) and \emph{perfectly balanced} (valence near zero)---a distinction the
system preserves by storing salience and valence \emph{separately}, so a consumer can tell silence
from contested framing. The momentum $\dot P$ is positive when a narrative is building and negative
when a backlash is forming. Because salience, valence, perception, and momentum are each persisted on
a fixed cadence, any of the four is independently available as a series.

\begin{remark}[Why the product]
\label{rem:product}
Coupling attention and tone \emph{multiplicatively} is the key modeling choice. An additive
combination would let high attention mask neutral tone, or strong tone with negligible coverage
register as a large signal. The product is zero whenever \emph{either} factor is zero---no coverage,
or perfectly balanced coverage---and is large only when a concept is both heavily covered and
directionally framed. This is the formal content of ``pressure.''
\end{remark}

% ============================================================
\section{Properties}
\label{sec:properties}
% ============================================================

\begin{proposition}[Single-source immunity]
\label{prop:single}
Both measures are immune to single-source amplification. In \eqref{eq:resonance} and
\eqref{eq:salience} the breadth factor $\log_2(1+n_{\text{sources}})$ multiplies the entire
contribution, so a cluster drawn from one source contributes the minimum breadth regardless of its
volume; and in \eqref{eq:valence} a single source's sentiment is down-weighted by its (necessarily
low, un-corroborated) trust from \eqref{eq:probation}.
\end{proposition}

\begin{proposition}[Full decomposability]
\label{prop:decomp}
Every measure is a composition of independently persisted primitives (breadth, decay, trust,
engagement, sentiment, salience, valence). Hence any reported value is reconstructable from stored
components, and the contribution of any factor is recoverable by ablation---the measure is auditable
by construction, never a black box.
\end{proposition}

These properties operationalize the platform's editorial stance as quantitative constraints: bias is
shown rather than hidden (decomposability), and no single voice can manufacture consensus (single-source
immunity).

% ============================================================
\section{Application: An Exogenous Narrative Signal for Markets}
\label{sec:application}
% ============================================================

Perception Pressure is, by construction, a per-concept float-valued time-series aligned to wall-clock
time. This makes it a candidate \emph{exogenous} predictor for any time-series problem whose entities
can be mapped to tracked concepts---most naturally, financial markets, where the question is whether
narrative leads price.

The intended use, realized in a companion quantitative-research system, is to align $P(K,t)$ (and its
components $A$, $V$, and momentum $\dot P$) to market bars for the matching entity by an as-of join,
and then to test each series for predictive content exactly as one would test an endogenous market
feature: by its information coefficient across forecast horizons and by its mutual information with
forward returns, under planted-signal and no-look-ahead controls. The narrative series is thus not
privileged; it must earn its place by the same gates as any other candidate signal.

\begin{remark}[Leading vs.\ contemporaneous]
\label{rem:leading}
The momentum term $\dot P$ is the natural \emph{leading} candidate: a narrative that is building
($\dot P > 0$) before it is fully priced is the economically interesting case, whereas a level $P$
that merely co-moves with price is contemporaneous and of less predictive use. Distinguishing the two
is a central question for the evaluation below.
\end{remark}

% ============================================================
\section{Evaluation Protocol}
\label{sec:eval}
% ============================================================

The repository implements the measures but contains no validation artifacts; we therefore specify the
evaluation rather than report it. Four studies would substantiate the measure:

\begin{enumerate}
  \item \textbf{Construct validity.} Trace $P(K,t)$ across known narrative episodes (a story building
        then fading; a sign flip on backlash) and confirm the trajectory matches the qualitative
        account---a sign flip in $V$, a rise then decay in $A$.
  \item \textbf{Decomposition.} Plot $A$, $V$, and $P$ jointly to show that the product separates
        ``uncovered'' from ``balanced,'' which a scalar tone index cannot.
  \item \textbf{Ablation.} Every constant (the half-lives, the engagement normalizer, the breadth
        base, the trust thresholds) is a configuration parameter; vary each and measure sensitivity of
        the resulting series, isolating the contribution of the breadth multiplier and the
        trust-weighting.
  \item \textbf{Predictive utility.} For matched market entities, test whether $\dot P$ carries
        information \emph{incremental} to price and volume for forward returns, under the
        information-coefficient and mutual-information gates of Section~\ref{sec:application}, with
        false-discovery-rate control across concepts and horizons.
\end{enumerate}

Study~4 is the strongest possible empirical claim and the one the companion system is built to run; it
is deliberately framed as leakage-controlled and multiple-testing-corrected, because a narrative
signal tested against many entities and horizons is otherwise a textbook data-snooping trap.

% ============================================================
\section{Limitations}
\label{sec:limitations}
% ============================================================

We disclose the principal threats to validity. Sentiment and bias labels are elicited from a large
language model, not from gold-standard human annotation, so they carry measurement error whose
structure is itself a research question. Story clustering and deduplication are lexical (token
overlap, TF--IDF, regular-expression entity extraction) rather than embedding-based, which can split
or merge events. The engagement term defaults to neutral when reaction counts are unavailable, so it
is often inert in practice. Finally, the implementation is single-writer at the storage layer, which
bounds throughput. None of these undermines the measures' definitions, but each bears on the fidelity
of the computed series and must be controlled in any empirical study.

% ============================================================
\section{Conclusion}
\label{sec:conclusion}
% ============================================================

We have defined two decomposable, trust-weighted measures of media narrative and described the system
that computes them continuously from a multi-source news stream. Perception Pressure couples weighted
attention with credibility-weighted directional sentiment into a single signed, magnitude-meaningful,
auditable signal; Resonance scores per-story impact with diversity as a precondition. Credibility is
earned through cross-corroboration rather than assigned. The measures are fully specified and
parameterized for ablation, and we have set out the construct-validity, decomposition, sensitivity,
and predictive-utility studies that would validate them---the last, a leakage-controlled test of
whether narrative momentum leads returns, being the sharpest empirical question. The contribution
here is the instrument and its methodology; turning the instrument on data is the next step.

% ============================================================
\section*{References}
\addcontentsline{toc}{section}{References}
% ============================================================

\begin{enumerate}
  \item Antweiler, W.\ \& Frank, M.Z.\ (2004). Is all that talk just noise? The information content of
        internet stock message boards. \textit{Journal of Finance}, 59(3), 1259--1294.

  \item Baker, S.R., Bloom, N.\ \& Davis, S.J.\ (2016). Measuring economic policy uncertainty.
        \textit{Quarterly Journal of Economics}, 131(4), 1593--1636.

  \item Benjamini, Y.\ \& Hochberg, Y.\ (1995). Controlling the false discovery rate: a practical and
        powerful approach to multiple testing. \textit{Journal of the Royal Statistical Society,
        Series B}, 57(1), 289--300.

  \item Bollen, J., Mao, H.\ \& Zeng, X.\ (2011). Twitter mood predicts the stock market.
        \textit{Journal of Computational Science}, 2(1), 1--8.

  \item Choi, H.\ \& Varian, H.\ (2012). Predicting the present with Google Trends.
        \textit{Economic Record}, 88(s1), 2--9.

  \item Da, Z., Engelberg, J.\ \& Gao, P.\ (2011). In search of attention.
        \textit{Journal of Finance}, 66(5), 1461--1499.

  \item Gentzkow, M., Shapiro, J.M.\ \& Taddy, M.\ (2019). Measuring group differences in
        high-dimensional choices: method and application to congressional speech.
        \textit{Econometrica}, 87(4), 1307--1340.

  \item Kraskov, A., St\"ogbauer, H.\ \& Grassberger, P.\ (2004). Estimating mutual information.
        \textit{Physical Review E}, 69(6), 066138.

  \item Loughran, T.\ \& McDonald, B.\ (2011). When is a liability not a liability? Textual analysis,
        dictionaries, and 10-Ks. \textit{Journal of Finance}, 66(1), 35--65.

  \item McCombs, M.E.\ \& Shaw, D.L.\ (1972). The agenda-setting function of mass media.
        \textit{Public Opinion Quarterly}, 36(2), 176--187.

  \item Preis, T., Moat, H.S.\ \& Stanley, H.E.\ (2013). Quantifying trading behavior in financial
        markets using Google Trends. \textit{Scientific Reports}, 3, 1684.

  \item Tetlock, P.C.\ (2007). Giving content to investor sentiment: the role of media in the stock
        market. \textit{Journal of Finance}, 62(3), 1139--1168.
\end{enumerate}

\end{document}
